\section{Proofs for the classical lower bound and the random-graph speedup}
\label{app:classical-lower}

\subsection{A weight bound for adaptive queries to a random graph}
\label{app:weight-bound}

We record the random-Painter estimate of Conlon, Fox, Grinshpun, and He
\cite[proof of Theorem~1]{ConlonFoxGrinshpunHe2019} in the form used in
this paper.  We follow their weight argument, taking the matching number
of the queried graph as the potential; removing the critical edge from a
maximum matching then leaves a matching of the smaller set, which keeps the
recursion closed.

Throughout this subsection $V\ge2$, and $G\sim G(V,1/2)$ assigns
independent uniform colours in $\{0,1\}$ to the pairs of $[V]$.  A
deterministic algorithm queries $T$ distinct pairs adaptively.  Let
$e_1,\ldots,e_T$ be the queried pairs, let $Q_t$ be the graph formed by
$e_1,\ldots,e_t$, and let $\mathcal F_t$ be the $\sigma$-algebra generated
by the first $t$ answers.  Because the algorithm is deterministic, $e_{t+1}$
is $\mathcal F_t$-measurable.

For $U\subseteq[V]$ and $0\le t\le T$, let $e_t(U)$ be the number of pairs
inside $U$ among $e_1,\ldots,e_t$, and call $U$ \emph{consistent at time
$t$} if all of these pairs have colour~$1$.  Define
\[
 w(U,t)=
 \begin{cases}
 2^{-\binom{|U|}2+e_t(U)}, & U\text{ is consistent at time }t,\\
 0,&\text{otherwise},
 \end{cases}
\]
the conditional probability, given $\mathcal F_t$, that $U$ becomes a
colour-$1$ clique when its remaining pairs are revealed.  Let $\nu(U,t)$ be
the matching number of $Q_t[U]$; it is nondecreasing in $t$ and increases by
at most one when a pair is added.  For integers $k\ge2c\ge0$ set
\[
 w_{k,c}(t)=\sum_{|U|=k,\ \nu(U,t)\ge c}w(U,t).
\]

\begin{lemma}[Weight bound]\label{lem:weight-bound}
For all integers $k\ge2c\ge0$,
\[
 \mathbb E\,w_{k,c}(T)\le 2^{-\binom k2+c(c-1)}\,V^{k-2c}\,T^{c}.
\]
\end{lemma}

\begin{proof}
\emph{Martingale property.}  Fix $U$.  If $e_{t+1}\not\subseteq U$, then
$w(U,t+1)=w(U,t)$.  If $e_{t+1}\subseteq U$, then with probability $1/2$
the new pair has colour~$1$ and $w(U,t+1)=2w(U,t)$, and otherwise
$w(U,t+1)=0$.  Hence $\mathbb E[w(U,t+1)\mid\mathcal F_t]=w(U,t)$, and in
particular $\mathbb E\,w(U,T)=w(U,0)=2^{-\binom{|U|}2}$.  Summing over all
$k$-sets gives the case $c=0$:
\[
 \mathbb E\,w_{k,0}(T)=\binom Vk2^{-\binom k2}\le V^k2^{-\binom k2}.
\]

\emph{Monotonicity.}  Since $\nu(U,\cdot)$ is nondecreasing,
$w_{k,c}(t+1)\ge\sum_{\nu(U,t)\ge c}w(U,t+1)$, and each event
$\{\nu(U,t)\ge c\}$ is $\mathcal F_t$-measurable.  Taking conditional
expectations gives $\mathbb E\,w_{k,c}(t)\le\mathbb E\,w_{k,c}(t+1)$.

\emph{Critical times.}  Let $c\ge1$.  For each $U$ let
$\tau(U)=\min\{t:\nu(U,t)\ge c\}$, a stopping time with $\tau(U)\ge1$.
Since $\{\tau(U)=t\}\in\mathcal F_t$ and $w(U,\cdot)$ is a martingale,
\[
 \mathbb E\,w_{k,c}(T)
 =\sum_{|U|=k}\sum_{t=1}^T\mathbb E\bigl[w(U,T)\,\mathbf 1[\tau(U)=t]\bigr]
 =\sum_{t=1}^T\mathbb E\,w^*_{k,c}(t),
\]
where
\[
 w^*_{k,c}(t)=\sum_{|U|=k,\ \tau(U)=t}w(U,t).
\]

\emph{Charging.}  Fix $t$ and write $e_t=\{x,y\}$.  Suppose $|U|=k$ and
$\tau(U)=t$.  Then $\nu(U,t-1)=c-1$, $\nu(U,t)=c$, and $e_t\subseteq U$.
Moreover, every maximum matching of $Q_t[U]$ contains $e_t$, since a
maximum matching avoiding $e_t$ would be a matching of size $c$ in
$Q_{t-1}[U]$.  Fix such a matching $M$, let $M'=M\setminus\{e_t\}$, let $A$
be the set of $2(c-1)$ vertices covered by $M'$, and let
$U'=U\setminus\{x,y\}$.  Then $M'$ is a matching of $Q_{t-1}[U']$, so
$\nu(U',t-1)\ge c-1$.  We claim that at most $2(c-1)$ of the pairs between
$\{x,y\}$ and $U'$ belong to $Q_{t-1}$.  Indeed, if $xa\in Q_{t-1}$ with
$a\in U'\setminus A$, then $M'\cup\{xa\}$ is a matching of size $c$ in
$Q_{t-1}[U]$, a contradiction; the same holds for $y$.  For each
$ab\in M'$, the pairs $xa$ and $yb$ cannot both lie in $Q_{t-1}$, since
$(M'\setminus\{ab\})\cup\{xa,yb\}$ would again be a matching of size $c$ in
$Q_{t-1}[U]$; likewise $xb$ and $ya$ cannot both lie in $Q_{t-1}$.  Hence
at most two of $xa,xb,ya,yb$ lie in $Q_{t-1}$, and the claim follows by
summing over the $c-1$ edges of $M'$.

Consequently $e_t(U)\le e_{t-1}(U')+1+2(c-1)$, and $U$ is consistent at
time $t$ only if $U'$ is consistent at time $t-1$ and $e_t$ has colour~$1$.
Using $\binom k2-\binom{k-2}2=2k-3$, we obtain
\[
 w(U,t)\le 2^{-(2k-2c-2)}\,w(U',t-1)
 \quad\text{if }e_t\text{ has colour }1,
\]
and $w(U,t)=0$ otherwise.  The map $U\mapsto U'$ is injective because
$U=U'\cup e_t$.  Therefore
\[
 w^*_{k,c}(t)\le
 \mathbf 1[e_t\text{ has colour }1]\;2^{-(2k-2c-2)}\,w_{k-2,c-1}(t-1).
\]
The indicator is independent of $\mathcal F_{t-1}$ with mean $1/2$, while
$w_{k-2,c-1}(t-1)$ is $\mathcal F_{t-1}$-measurable.  Taking expectations
and using monotonicity,
\[
 \mathbb E\,w^*_{k,c}(t)\le 2^{-(2k-2c-1)}\,\mathbb E\,w_{k-2,c-1}(T),
\]
and summing over $t\le T$ gives
\[
 \mathbb E\,w_{k,c}(T)\le T\,2^{-(2k-2c-1)}\,\mathbb E\,w_{k-2,c-1}(T).
\]

\emph{Iteration.}  Applying the last inequality $c$ times, the exponents
$2(k-2i)-2(c-i)-1$ for $i=0,\ldots,c-1$ sum to $2kc-3c^2$, and the case
$c=0$ gives
\[
 \mathbb E\,w_{k,c}(T)
 \le T^c\,2^{-(2kc-3c^2)}\,V^{k-2c}\,2^{-\binom{k-2c}2}
 =2^{-\binom k2+c(c-1)}\,V^{k-2c}\,T^c,
\]
because $2kc-3c^2+\binom{k-2c}2=\binom k2-c(c-1)$.
\end{proof}

\begin{corollary}[Fully exposed monochromatic sets]\label{cor:exposed}
Let $K\ge2$, let $V\ge K$, and let a deterministic algorithm make $T\ge1$
adaptive edge queries to $G\sim G(V,1/2)$.  For every integer
$0\le c\le K/2$, the expected number of $K$-sets all of whose $\binom K2$
pairs have been queried and carry the same colour is at most
\[
 2\cdot2^{-\binom K2+c(c-1)}(2T)^{K-c}.
\]
\end{corollary}

\begin{proof}
Renaming vertices in order of first appearance in a query does not change
the joint distribution of the queried graph and its colours, because the
colours of distinct pairs are independent and uniform; the renamed
algorithm queries pairs inside $[V']$ with $V'=\min\{V,2T\}$, and the
number of fully exposed monochromatic $K$-sets is unchanged.  A fully
exposed colour-$1$ $K$-set $U$ satisfies $\nu(U,T)=\lfloor K/2\rfloor\ge c$
and $w(U,T)=1$, so the number of such sets is at most $w_{K,c}(T)$.
\Cref{lem:weight-bound} with $V'\le2T$ bounds its expectation by
$2^{-\binom K2+c(c-1)}(2T)^{K-2c}T^c\le2^{-\binom K2+c(c-1)}(2T)^{K-c}$.
The same bound holds for colour~$0$ by symmetry.
\end{proof}

\subsection{Proof of the randomized lower bound}
\label{app:classical-lower-proof}

\begin{proof}[Proof of \cref{prop:classical-lower}]
Let $\beta=2-\sqrt2$ and $\alpha=1-1/\sqrt2$, so that $\beta=2\alpha$ and
$\alpha^2-2\alpha+1/2=0$.  Set
\[
 T=\bigl\lfloor2^{\beta K-2}\bigr\rfloor,
 \qquad
 c=\lfloor\alpha K\rfloor=\alpha K+\epsilon,
 \qquad -1<\epsilon\le0.
\]
Since $\alpha<1/2$, we have $0\le c\le K/2$, and $K\ge8$ gives $T\ge1$.

Consider first a deterministic algorithm making at most $T$ queries to
$G\sim G(N,1/2)$; padding with unused queries, we may assume that it makes
exactly $T$.  By \cref{cor:exposed} and Markov's inequality, the
probability that the queried pairs contain a fully exposed monochromatic
$K$-set is at most
\begin{equation}\label{eq:random-painter-bound}
 2\cdot
 2^{-\binom K2+c(c-1)}(2T)^{K-c}.
\end{equation}
Since $2T\le2^{\beta K-1}$, the base-$2$ logarithm of
\cref{eq:random-painter-bound} is at most
\[
 1-\binom K2+c(c-1)+(K-c)(\beta K-1)
 =1-\frac K2+\epsilon^2,
\]
where the identity follows by substituting $c=\alpha K+\epsilon$ and
$\beta=2\alpha$ and using $\alpha^2-2\alpha+1/2=0$.  For $K\ge8$ this
quantity is at most $-2$, so the event has probability at most $1/4$.

Now condition on the query transcript and suppose that the algorithm
outputs a $K$-set whose internal pairs have not all been exposed in one
colour.  If the exposed internal pairs contain both colours, the output is
not homogeneous.  If they contain exactly one colour, at least one
unexposed pair must match that colour, which has conditional probability at
most $1/2$, because unexposed pairs are independent fair coins given the
transcript.  If no internal pair has been exposed, the success probability
is $2^{1-\binom K2}\le1/2$ for $K\ge3$.  Thus the conditional success
probability is at most $1/2$ outside the fully exposed event, and the
success probability under $G(N,1/2)$ is at most
\[
 \frac14+\frac34\cdot\frac12=\frac58.
\]

A randomized algorithm making at most $T$ queries is a distribution over
such deterministic algorithms, so its success probability on
$G\sim G(N,1/2)$ is also at most $5/8$.  Finally, an algorithm whose
worst-case success probability on graphs with $M$ vertices is at least
$2/3>5/8$ must make more than $T$ queries, and $M=4^{K-1}$ gives
\[
 2^{\beta K}
 =\Theta\!\left(M^{\beta/2}\right)
 =\Theta\!\left(M^{1-1/\sqrt2}\right),
\]
which proves the proposition.
\end{proof}

\subsection{Proof of the random-graph speedup}
\label{app:random-graph}

\begin{proof}[Proof of \cref{prop:random-graph}]
Let $M=4^{K-1}$.  For $W\subseteq[M]$ let
\[
 C(W)=\{u\in[M]\setminus W:\ G(w,u)=1\text{ for every }w\in W\}
\]
be the common neighbourhood of $W$.  Call $G$ \emph{expanding} if
$|C(W)|\ge M2^{-j-1}$ for every $j\in\{1,\ldots,K-1\}$ and every $j$-set
$W$; for $W=\emptyset$ the bound $|C(\emptyset)|=M$ holds trivially.

\emph{Random graphs are expanding.}  Fix $1\le j\le K-1$ and a $j$-set
$W$.  Then $|C(W)|\sim\operatorname{Bin}(M-j,2^{-j})$ with mean
$\mu=(M-j)2^{-j}$.  Since $M=4^{K-1}\ge4(K-1)\ge4j$, we have
$\mu\ge\frac34M2^{-j}$, hence $M2^{-j-1}\le\frac23\mu$ and
$\mu\ge\frac34M2^{-(K-1)}=\frac32\cdot2^{K-2}$.  The Chernoff bound
$\Pr[X\le(1-\gamma)\mu]\le e^{-\gamma^2\mu/2}$ \cite{Chernoff1952} with
$\gamma=1/3$ gives
\[
 \Pr\bigl[|C(W)|<M2^{-j-1}\bigr]\le e^{-\mu/18}\le e^{-2^{K-2}/12}.
\]
A union bound over the at most $\sum_{j=1}^{K-1}M^j\le2M^{K-1}$ sets $W$
shows that $G$ fails to be expanding with probability at most
$\pi_K=2M^{K-1}e^{-2^{K-2}/12}$.

\emph{The algorithm.}  Set $W_0=\emptyset$.  For $j=0,\ldots,K-1$, given
the clique $W_j=\{v_1,\ldots,v_j\}$, run the capped search of
\cref{lem:capped-search} on $[M]$ for the predicate ``$u\notin W_j$ and
$G(v_i,u)=1$ for all $i\le j$'', with density parameter
$\lambda_j=2^{-j-1}$, repeating the block $O(\log(K/\eta))$ times so that
its failure probability is at most $\eta/K$ whenever the density promise
holds.  The predicate costs $O(j+1)$ edge queries, including uncomputation.
Let $v_{j+1}$ be the verified output and $W_{j+1}=W_j\cup\{v_{j+1}\}$.  If
any search fails, output $\bot$; otherwise output $W_K$.

\emph{Correctness and cost.}  Every accepted vertex was verified to be
adjacent to all earlier vertices, so a non-$\bot$ output is a $K$-clique
for every input graph.  If $G$ is expanding, then the search at level $j$
has marked density at least $\lambda_j$, and a union bound over the $K$
searches shows that the algorithm outputs $W_K$ with probability at least
$1-\eta$.  Hence it succeeds with probability at least $1-\eta-\pi_K$ over
$G$ and its internal randomness.  Each capped block at level $j$ uses
$O(2^{j/2})$ predicate evaluations, so the total number of edge queries is
\[
 O\!\left(\log\frac K\eta\sum_{j=0}^{K-1}(j+1)2^{j/2}\right)
 =O\!\left(K2^{K/2}\log\frac K\eta\right)
\]
on every input, because every block has a predetermined cap.
\end{proof}
